%% 26장 희소 행렬과 직접법
\chapter{희소 행렬과 직접법}
\label{ch:sparse}

\abstract{편미분방정식을 이산화한 행렬, 그래프의 인접 행렬, 큰 회귀 문제의 설계 행렬은 대부분의 원소가 $0$이다. 이런 \term{희소 행렬}{sparse matrix}을 밀집 행렬처럼 저장하고 계산하면 메모리와 시간을 $0$에 낭비한다. 이 장에서는 $0$이 아닌 원소만 저장하는 형식들과 그 위의 기본 연산을 보고 희소 행렬을 LU나 Cholesky로 분해할 때 생기는 가장 큰 문제인 채움을 그래프의 언어로 설명한다. 같은 행렬도 행과 열의 순서를 바꾸면 채움이 수백 배 달라진다는 것을 직접 계산으로 보이고 채움을 줄이는 순서 매기기인 최소 차수와 중첩 분할을 다룬다. 마지막으로 희소 LU에서 안정성과 희소성을 저울질하는 문턱 피벗, 분해의 결과를 믿을 수 있는지 판정하는 조건수 추정과 반복 다듬기를 stardust가 Apple Accelerate를 감싸는 방식과 함께 본다.}

\section{저장 형식}
\label{sec:sparse-formats}

\subsection{COO, CSR, CSC}

가장 단순한 형식은 $0$이 아닌 원소마다 (행, 열, 값)의 세 쌍을 늘어놓는 \term{좌표 형식}{coordinate format, COO}이다. 만들기는 쉽지만 같은 행의 원소를 찾으려면 전체를 훑어야 한다. 계산에는 \term{압축 행 형식}{compressed sparse row, CSR}이 표준이다. 원소들을 행 순서로 늘어놓고 각 원소의 열 번호 배열 \code{indices}와 값 배열 \code{data}, 그리고 행 $i$의 원소가 \code{indices}의 어디서 시작하는지를 담은 길이 $m+1$의 배열 \code{indptr}을 둔다. 행 $i$의 원소는 구간 $[\code{indptr}_i,\code{indptr}_{i+1})$에 있다. 12장의 문자열 위치표와 같은 구조이다. 열을 기준으로 같은 일을 하면 \term{압축 열 형식}{compressed sparse column, CSC}이고 CSR 행렬의 세 배열은 그 전치 행렬의 CSC 배열과 같다.

COO에서 CSR로 바꾸는 일은 15장의 계수 정렬이다. 행마다 원소의 개수를 세어 누적하면 \code{indptr}이 되고 원소를 그 자리에 흩뿌리면 된다. 계수 정렬은 안정적이므로 같은 행 안의 순서는 입력 순서로 남는다. 각 행 안에서 열 번호가 오름차순이고 같은 자리의 원소가 하나뿐인 꼴을 \term{표준 꼴}{canonical form}이라 하는데, 행 안을 정렬하고 같은 자리의 중복을 더하면 얻는다. 표준 꼴을 유지하면 행 안에서 이분 탐색으로 원소를 찾을 수 있고 두 행렬을 병합하는 연산이 단순해지며 다음 절에서 보듯 합하는 순서가 정해져 결과가 결정적이 된다.

\subsection{색인의 폭}

원소 하나를 저장하는 데는 값 $8$바이트와 색인이 든다. 색인을 $32$비트 정수로 두면 원소당 $12$바이트이고 $64$비트로 두면 $16$바이트이다. 다음 절의 행렬--벡터 곱처럼 메모리 대역에 묶인 연산에서는 이 차이가 속도에 그대로 나타날 수 있다. 반면 $32$비트 색인은 원소가 $2^{31}$개를 넘는 행렬을 담지 못한다.

\begin{instardust}[\code{docs/design/sparse.md} §2.1--§2.3]
stardust의 희소 영역은 CSR, CSC, COO, DOK 형식을 두고 CSR과 CSC는 언제나 표준 꼴이다. scipy는 검사 없이 만든 비표준 꼴을 두었다가 연산마다 필요하면 정리하지만 stardust는 생성자에서 범위를 검사하고 정렬되지 않았거나 중복이 있으면 정리한다. 하나씩 넣어 만드는 DOK는 16장의 선형 탐사 해시 테이블을 C++에 둔다. 색인은 언제나 \code{int64}인데, 이 라이브러리의 정수 자료형이 \code{int64} 하나뿐이기 때문이다. 그 대가를 설계 문서는 측정으로 적어 둔다. 같은 행렬--벡터 곱 커널을 두 폭으로 재면 한 스레드에서 $1.02$--$1.16$배, 열 스레드에서는 $1.33$--$1.39$배 느렸다.
\end{instardust}

\section{희소 연산}
\label{sec:sparse-ops}

CSR 행렬과 밀집 벡터의 곱 $y=Ax$는 행마다 저장 위치 $p$에 대해 $y_i=\sum_p\code{data}_p\,x_{\code{indices}_p}$를 계산한다. 원소마다 값과 색인 하나씩, 그리고 $x$의 흩어진 원소 하나를 읽고 곱셈 하나와 덧셈 하나를 하므로, 13장의 언어로 이 연산은 메모리 대역에 묶여 있다. 행들은 서로 독립이므로 행을 나누어 여러 스레드에서 계산할 수 있고 행 안의 합을 표준 꼴의 열 순서로 차례로 더하면 결과는 스레드 수와 무관하다(14장).

두 희소 행렬의 곱 $C=AB$를 CSR로 계산하는 표준적인 방법은 Gustavson의 알고리즘이다\cite{gustavson1978}. $C$의 $i$번째 행은 $A$의 $i$번째 행의 원소 $a_{ik}$마다 $B$의 $k$번째 행에 $a_{ik}$를 곱해 더한 것이다. 길이 $n$의 밀집 누산 배열과, 이번 행에서 어느 열이 이미 채워졌는지를 적는 표지 배열을 두면, 행 하나를 $B$의 해당 행들의 원소 수에 비례하는 시간에 계산할 수 있다. 결과 행의 열 번호는 마지막에 정렬해 표준 꼴로 만든다. 전치는 CSR을 CSC로 다시 읽는 것이므로 계수 정렬 한 번이다.

\section{채움}
\label{sec:fill-in}

\subsection{소거와 그래프}

희소 행렬을 19장의 LU나 20장의 Cholesky로 분해하면, 원래 $0$이던 자리가 소거하는 동안 $0$이 아니게 될 수 있다. 이것을 \term{채움}{fill-in}이라 한다. 대칭 행렬로 생각하면 채움을 그래프로 정확히 설명할 수 있다\cite{parter1961}. 행렬 $A$에 대해 정점 $1,\dots,n$을 두고 $a_{ij}\ne0$이면 $i$와 $j$를 잇는 그래프를 만들자. 정점 $k$를 소거하면, 곧 $k$번째 피벗으로 남은 행렬을 갱신하면 $a_{ij}\gets a_{ij}-a_{ik}a_{kj}/a_{kk}$이므로, $k$의 이웃인 $i$와 $j$ 사이에 원래 간선이 없었더라도 간선이 생긴다. 곧 소거는 정점 하나를 지우면서 그 이웃들을 모두 서로 잇는 일이고 새로 생긴 간선이 채움이다. 분해된 인수 $L$의 $0$이 아닌 원소는 원래의 간선과 채움 간선을 합친 것이다.

\subsection{순서가 결정한다}

소거하는 순서는 결과의 채움을 크게 바꾼다. 극단적인 예가 정점 하나가 다른 모든 정점과 이어진 \term{화살 행렬}{arrow matrix}이다. 중심 정점을 맨 먼저 소거하면 나머지 모든 정점이 서로 이어져 $L$이 꽉 찬다. $n=1000$이면 $L$의 원소가 $500500$개이다. 중심 정점을 맨 나중에 소거하면 각 정점의 이웃이 중심 하나뿐이어서 채움이 전혀 없고, $L$의 원소는 $1999$개이다. 같은 행렬, 같은 분해인데 순서만으로 $250$배 차이가 난다.

현실적인 예를 보자. $30\times30$ 격자 위의 2차원 Laplace 연산자를 이산화한 행렬($n=900$)은 아래 삼각의 원소가 $2640$개이다. 행 순서대로 소거하면 $L$의 원소가 $27029$개로 늘지만 다음 절의 최소 차수 순서로 소거하면 $10351$개이다. 이 수치들은 모두 직접 짠 기호적 소거로 센 것이다. 행 순서의 소거에서 $L$은 행렬의 띠 안을 거의 다 채우는데, 격자의 띠 폭이 $\sqrt n$이므로 채움이 $n^{3/2}$로 자라고 연산은 $n^2$으로 자란다.

\subsection{최소 차수와 중첩 분할}

채움을 가장 적게 하는 순서를 찾는 문제는 NP-어려운 문제로 알려져 있으므로 휴리스틱을 쓴다. 가장 오래되고 널리 쓰이는 것은 매 단계에서 이웃이 가장 적은 정점, 곧 \term{차수}{degree}가 가장 작은 정점을 소거하는 \term{최소 차수}{minimum degree} 순서이다\cite{tinney1967}. 차수가 작은 정점을 소거하면 새로 이어지는 이웃 쌍이 적기 때문이다. 정확한 차수를 매번 갱신하는 것은 비싸서 오늘날은 차수의 근삿값과 정점들을 묶어 다루는 기법으로 빠르게 만든 \term{근사 최소 차수}{approximate minimum degree, AMD}를 쓴다\cite{amestoy1996amd}.

다른 접근은 분할 정복이다. 그래프를 작은 \term{분리 집합}{separator}으로 두 부분으로 나누고 두 부분을 먼저 소거한 뒤 분리 집합을 마지막에 소거한다. 두 부분 사이에는 간선이 없으므로 채움이 각 부분 안과 분리 집합에만 생긴다. 이것을 재귀적으로 되풀이하는 \term{중첩 분할}{nested dissection}은 2차원 격자에서 채움을 $O(n\log n)$, 연산을 $O(n^{3/2})$로 줄이고 이것이 그 격자에 대해 점근적으로 최적이다\cite{george1973nested}. 일반적인 그래프의 분할은 Metis 같은 다층 그래프 분할로 한다\cite{karypis1998metis}.

분해의 실제 계산은 \term{제거 나무}{elimination tree}로 정리된다\cite{liu1990etree}. 열 $j$의 부모를 $L$에서 $j$ 아래의 첫 $0$이 아닌 원소의 행으로 두면 나무가 되는데, 이 나무에서 서로 조상 관계가 아닌 열들은 독립적으로 계산할 수 있다. 또 구조가 같은 이웃 열들을 \term{초마디}{supernode}로 묶으면 그 안의 계산이 밀집 행렬 연산이 되어, 19장의 3층 BLAS의 속도로 분해할 수 있다\cite{davis2006direct}.

\section{희소 LU의 실제}
\label{sec:sparse-lu}

\subsection{문턱 피벗}

비대칭 희소 행렬의 LU는 두 요구가 부딪힌다. 19장의 부분 피벗은 안정성을 위해 열에서 가장 큰 원소를 피벗으로 고르라고 하지만 채움을 줄이려면 미리 정한 순서에 가까운 원소를 골라야 한다. 절충안은 \term{문턱 부분 피벗}{threshold partial pivoting}이다. 열에서 가장 큰 원소의 $\tau$배 이상인 후보들 가운데 채움을 가장 적게 만드는 것을 고른다. $\tau=1$이면 부분 피벗이고, $\tau$가 작을수록 희소성을 더 지키는 대신 성장 인수의 상한이 $(1+1/\tau)$의 거듭제곱으로 커진다. 흔히 $\tau=0.1$이나 $0.01$을 쓴다.

\subsection{판정과 다듬기}

희소 분해에서 결과를 믿을 수 있는지를 판정하는 것은 밀집 분해보다 까다롭다. 피벗이 반올림 때문에 정확히 $0$이 되지 않는 정확히 특이한 행렬도 있고 문턱 피벗의 성장 때문에 후진 오차가 커지는 행렬도 있다. 25장의 1-노름 추정을 분해된 인수의 풀기 몇 번으로 돌려 $\kappa_1(A)$를 추정하면 특이성을 판정할 수 있고, 풀이의 \term{성분별 후진 오차}{componentwise backward error}
\[
\omega=\max_i\frac{|b-A\hat x|_i}{(|A||\hat x|+|b|)_i}
\]
를 계산하면 계산한 해가 원소마다 상대적으로 얼마만큼 바뀐 문제의 정확한 해인지를 알 수 있다\cite{oettli1964}. $\omega$가 단위 반올림의 몇 배를 넘으면 19장과 22장에서 본 반복 다듬기를 한 번 하면 거의 언제나 $u$ 수준으로 돌아온다.

\begin{instardust}[\code{docs/design/sparse.md} §1.4, §3.6]
stardust의 희소 직접 풀이는 Apple Accelerate의 희소 LU(문턱 피벗 $\tau=0.01$)를 AMD 순서로 부른다. 설계 문서에는 Accelerate의 순서 매기기들의 결정성을 잰 표가 있는데 같은 입력을 같은 프로세스에서 스무 번 분해했을 때 AMD 순서의 LU는 한 번도 다른 비트를 내지 않았지만 기본 순서인 다중 스레드 Metis는 매번 다른 비트를 냈고 여러 스레드에서 동시에 부르면 프로세스가 죽었다. 그래서 순서를 고르는 옵션 없이 AMD만 쓴다. 3차원 격자에서는 AMD의 채움이 Metis보다 많아 $1.5$--$2$배 느리다는 것이 그 대가이다. Accelerate는 수치적으로 특이한 행렬에서도 성공을 보고하므로, 분해 뒤에 \code{lacn2}로 $\kappa_1$을 추정해 $\kappa_1\varepsilon\ge1$이면 특이로 판정한다. 비용은 분해 시간의 $10$--$15\%$이다. 풀이마다 위의 성분별 후진 오차를 계산해 $8\varepsilon$를 넘으면 한 번 다듬는데, 설계 문서의 시험 행렬에서 안장점 행렬은 첫 풀이의 후진 오차가 $4362\varepsilon$였고 다듬은 뒤에는 모두 $0.6\varepsilon$ 이하였다.
\end{instardust}

\section{맺음}

희소 행렬의 계산은 행렬을 그래프로 보는 순간 명료해진다. 저장 형식은 그래프의 인접 목록이고 소거는 정점을 지우며 이웃을 잇는 일이며 좋은 순서는 그래프를 잘 나누는 일이다. 밀집 행렬의 분해에서 피벗이 안정성만의 문제였다면, 희소 행렬에서는 안정성과 희소성을 함께 저울질해야 한다. 그런데 분해 자체가 너무 비싸거나 채움이 너무 크면 직접법을 포기해야 하는데, 다음 장에서는 행렬--벡터 곱만으로 해에 다가가는 반복법을 다룬다.

희소 직접법은 Davis의 책\cite{davis2006direct}이 가장 친절하고 완전한 입문이다. 실제 응용에서 모은 희소 행렬의 방대한 모음도 같은 저자가 관리한다\cite{davis2011collection}.

\chapterreferences
